\begin{textsample}{POS Dim 8 – global\_south\_2025\_03 – Score 4.00 – global\_south\_2025\_03/122905981.txt}  \label{ex:f8_pos_147}
How reporting in Gaza is a deadly assignment for journalists Over the past year and a half, the world has watched the war in Gaza unfold.
It is a conflict that has killed more than 50,000 people, according to the Palestinian territory ' s health ministry, and almost all of the 2.3 million residents there have been displaced multiple times.
Local Palestinian journalists are the only ones capable of providing crucial \textbf{insight} on what is happening in Gaza to the world.
But the conflict, triggered by the attacks on Israel the Islamic militant group Hamas led on October 7, 2023, has made the small territory one of the most dangerous places in the world to report, according to the Committee to Protect Journalists ( CPJ ), a New York-based press freedom monitor.
CPJ said more than 170 journalists and media workers have been killed in Gaza since the beginning of the conflict.
Reporters Without Borders, another press freedom organization headquartered in Paris, puts the number at more than 200 journalists. @ @ @ @ @ @ @ @ @ @ reporters due to the lack of communication and electricity.
Some journalists have lost family members, friends and homes.
Internal pressures in Gaza, which is sealed off by Israel and Egypt, have added to the difficult environment in which journalists have to operate."It ' s hard to describe what it feels like to be in Gaza.
The constant noise of the bombing, the explosions, the number of people killed, it ' s indescribable,"Safinaz al Louh, a freelance journalist in Gaza, told DW before the recent ceasefire between Israel and Hamas was agreed in January.
She lost her brother, a cameraman, in the war.
Gaza Strip : One woman finds her voice in journalismTo view this video \textbf{please} enable JavaScript, and consider upgrading to a web browser that supports HTML5 video Another Palestinian journalist, Salma al Qaddoumi, told DW that the \textbf{displacement} and \textbf{separation} from family had been difficult to cope with during the 15 months of war before the pause in the fighting."@ @ @ @ @ @ @ @ @ @ one place, and then you have to start all over again, knowing that nowhere is really safe,"Al Qaddoumi said via WhatsApp from Gaza.
She was injured while reporting in southern Gaza during the war in an incident in which her colleague was killed.
As fighting resumes following the collapse of a ceasefire, the dangerous, sometimes deadly reporting conditions have returned.
Journalists accused of terrorism On March 24, two Palestinian journalists were killed in Gaza in two consecutive Israeli strikes.
One journalist, Mohammed Mansour was killed with his wife and son in an airstrike on his home in Khan Younis in southern Gaza.
Mansour worked for Palestine Today, a station affiliated with the Palestinian militant group Islamic Jihad.
Later in the day, Hossam Shabat, a 23-year-old correspondent for the Al Jazeera Mubasher channel, was killed in an airstrike on his car in Beit Lahiya, a city in northern Gaza.
In a statement, the Israel Defense Forces ( IDF ) said that it had"eliminated @ @ @ @ @ @ @ @ @ @ terrorist,"a charge he and Al Jazeera previously denied.
After Shabat ' s death, Al Jazeera, a Qatar-based news network banned in Israel and by the Palestinian Authority in the occupied West Bank, strongly condemned what it called an assassination of its correspondent.
It also called on the international community to condemn Israel ' s"systemic killing of journalists.""The international community must act fast to ensure that journalists are kept safe and hold Israel to account for the deaths of Hossam Shabat and Mohammed Mansour, whose killings may have been targeted.
Journalists are civilians, and it is illegal to attack them in a war zone,"the statement continued.
The CPJ called for an investigation into whether the IDF deliberately killed the journalists.
The non-profit organization released a report in February documenting at least 13 cases in Gaza and Lebanon where journalists were"deliberately targeted"by Israeli forces, adding that it was investigating more intentional killings.
The IDF rejected the CPJ accusations at the time in @ @ @ @ @ @ @ @ @ @ never, and will never, deliberately target journalists."The statement went on to say that the IDF"does not target civilian objects and civilians, including media organizations and journalists as such."Journalists in Gaza and press freedom organizations feel reporters are being unfairly labeled as terroristsImage : Dawoud Abo Alkas/Anadolu/picture alliance Several press freedom groups have questioned labeling journalists as"terrorists"and the implication this has on the safety of journalists."In many countries, there are journalists who operate as mouthpieces for authorities, political opposition or, in some cases, militant groups,"CPJ chief executive Jodie Ginsburg told DW at the end of January 2025, before the latest offensive."Unless they ' re engaged in direct incitement to violence or they ' re actually part of militant activity, that does not make them targets for killing."Ban on foreign journalists continues Local Palestinian journalists have been responsible for reporting on the war to a global audience.
That is because the Israeli government has maintained its ban @ @ @ @ @ @ @ @ @ @ press freedom organizations worldwide for unfettered access.
Israel ' s Supreme Court has yet to rule on a petition filed by the Foreign Press Association in Israel and the Palestinian territories demanding independent access for foreign media.
Until now, the Israeli army has only allowed some foreign and Israeli journalists into Gaza as part of military embedded visits, which are tightly controlled and do not allow journalists to more around independently."That level of restriction is totally unprecedented,"Ginsburg said."Certainly, when you talk to war correspondents that have covered everything from Chechnya to Sudan, not being able to have any access at all is completely unprecedented."Ginsburg said that because of this, all pressure rests on local Palestinian journalists to report what was happening."Because they are Palestinian and local journalists, they have this additional doubt cast on what they ' re reporting -- on top of which, of course, they ' re reporting on conditions of war."

% matched lemmas: displacement, insight, please, separation
\end{textsample}
